\section{Method}
\label{sec:method}

\paragraph{Probes.}
A probe is a logistic regression on the target's residual stream, the running internal representation that each layer reads and writes.
The last-prompt probe reads the activation at the last token of the model's input for that call, just before the model starts to write, so it scores a call before the call generates anything (\cref{eq:probe}).
We choose the layer inside each training split, using only training episodes, and weight leaking and non-leaking calls equally within each round \citep[following][]{goldowskydill2025}.
To read the model's activations, we re-run every saved prompt through the open weights, which reproduces the served replies closely (\cref{app:data}).

\paragraph{Reasoning probe.}
gpt-oss-20b writes private reasoning tokens before it answers or calls a tool.
On this model a second probe applies $w^\top h_\ell + b$ to every token of the call's reasoning $z_t$ and keeps the maximum.
Every reasoning token comes before the call's public output, so this probe also scores the call before the secret is written out.

\paragraph{Zero-shot text monitor.}
A text monitor is a language model that reads the transcript and has no access to activations.
Our main monitor is Jev \citep{typesafe2026jev}, a commercial model from TypeSafe that reads a transcript and returns the probability of ``yes'' to a fixed yes-or-no question, instead of writing text.
For each call, we show Jev the conversation and tool results that the target had seen, with the secret masked, and ask whether the agent's next reply will write or post the protected value.
On gpt-oss-20b, Jev also sees the target's earlier reasoning in the same turn, which the model's prompt retains.
It is used zero-shot, with no training on our labels.
\Cref{app:monitors} gives the prompts and two further monitors, an LLM judge and a second decision model.

\paragraph{Trained text classifier.}
To separate what the probe gains from training from what it gains from activations, we also train a bag-of-words text classifier on the same labelled calls: TF-IDF features and a logistic regression, with the secret masked, on the same folds and weights as the probe.
It reads either the call's full input or only the latest message of that input; the figures show the latest message on Qwen3-32B and the full input on gpt-oss-20b, and \cref{app:shortcuts} reports both views for both models.

\paragraph{A fair comparison.}
Probes, the monitor and the classifier are scored on exactly the same calls, with the same labels.
Leaks tend to come in early rounds, so a rule that simply rates early rounds as risky would score well.
We therefore compare each call only with calls from the same scenario and round.
We compute the AUROC within each such group $g$ and average the groups, weighting each by its number of positive calls $n^+_g$:
\begin{equation*}
  \mathrm{AUROC} = \sum_g \frac{n^+_g}{\sum_{g'} n^+_{g'}}\, \mathrm{AUROC}_g .
\end{equation*}
AUROC is the chance that a random positive call outranks a random negative one, so 0.5 is chance and 1 is perfect.
In words, the metric asks how often a leaking call outranks a non-leaking call from the same scenario and round.
Groups with only leaking or only non-leaking calls are dropped (\cref{tab:group-coverage} counts what remains), and any score that depends only on the scenario and the round therefore gets exactly 0.5.

\paragraph{Held-out scenarios.}
We train on five scenarios, test on the sixth, and rotate through all six, so every test call comes from a scenario that the probe never saw in training.
Related episodes, such as variants of one seeded episode or reruns that share a history, are always kept on the same side of a split.
Intervals are 95\% bootstrap intervals that resample these groups of related episodes (\cref{app:probes}).

\paragraph{Evidence.}
Test episodes are scored once, by \emph{frozen} detectors whose weights and alarm thresholds were fixed on development episodes beforehand.
The gpt-oss-20b reasoning test was \emph{pre-registered}: a committed plan fixed its analysis and success criterion before its episodes were collected (\cref{app:confirmation}).
In every figure, filled markers are test episodes and hollow markers are development episodes.
